\scanpage{47}
\diagram{\resizebox{\linewidth}{!}{\begin{tikzcd}[ampersand replacement=\&,column sep=small,row sep=1.5em]
\& 0 \& 0 \& 0 \& 0\\
0\arrow[r] \& \Phi(X)/(\mathrm G_m(X)/U)\arrow[u]\arrow[r] \& \Gamma(X,\underline{\Phi}_X/(O_X^*/U_X))\arrow[u]\arrow[r] \& \operatorname{Pic}(X)\arrow[u]\arrow[r,"{\deleted{\unsure{0\to H^1|\longrightarrow}}}" description] \& H^1(X,\underline{\Phi}_X)\arrow[u]\\
0\arrow[r,dash] \& \Phi(X)\arrow[u]\arrow[r] \& \mathrm{QF}(X)\arrow[u]\arrow[r] \& \operatorname{Div}(X)\arrow[u]\arrow[r] \& H^1(X,\underline{\Phi}_X)\arrow[u,phantom,"{=}" description,sloped]\\
0\arrow[r] \& \mathrm G_m(X)/U\arrow[u]\arrow[r] \& R(X)^*/U\arrow[u]\arrow[r] \& \operatorname{Div}_{\ell}(X)\arrow[u]\arrow[r] \& 0\arrow[u]\\
\& 0\arrow[u] \& 0\arrow[u] \& 0\arrow[u] \&
\end{tikzcd}}}

(utilisant que $\underset{\substack{\parallel\\\operatorname{Pic}(X)}}{H^1(X,\underline O_X^{\times})}\simeq H^1(X,\underline O_X^{\times}/U)$ grâce au fait que $U$ est flasque).

Je dis que \uline{si $X$ est quasi-projective}, alors dans le diagramme comm. précédent on peut remplacer $H^1(X,\underline{\Phi}_X)$ par 0 (deux fois), en d'autres termes que $\Gamma(X,\underline{\Phi}_X/(\underline O_X^{\times}/U_X))\to\operatorname{Pic}(X)$ est surjectif. Comme tt $\unsure{\xi|\zeta}\in\operatorname{Pic}(X)$ est la différence de deux éléments très amples, on est ramené à prouver que tt élément très ample est dans l'image, et
% archive sheet number: 46
